% Generated draft: requires mathematical and source review.
\documentclass[11pt]{article}
\usepackage[margin=25mm]{geometry}
\usepackage{fontspec}
\setmainfont{Latin Modern Roman}
\usepackage{amsmath,amssymb,mathtools,unicode-math}
\setmathfont{Latin Modern Math}
\usepackage{enumitem,needspace,xurl,hyperref}
\setlength{\parindent}{0pt}
\setlength{\parskip}{5pt}
\setlength{\emergencystretch}{3em}
\newcommand{\N}{\mathbb{N}}
\newcommand{\Z}{\mathbb{Z}}
\newcommand{\Q}{\mathbb{Q}}
\newcommand{\R}{\mathbb{R}}
\newcommand{\C}{\mathbb{C}}
\newcommand{\F}{\mathbb{F}}
\newcommand{\A}{\mathbb{A}}
\newcommand{\PP}{\mathbb{P}}
\newcommand{\E}{\mathbb{E}}
\newcommand{\eps}{\varepsilon}
\newcommand{\dd}{\,\mathrm{d}}
\begin{document}
\section*{UnsolvedMath ID: 10400155}
\textbf{Source problem code:} catalogue entry
\subsection{Definitions and mathematical statement}
Let \(H_k=H^0(\mathcal M_\Sigma,\mathcal L^{\otimes k})\), the space of global sections of the source's line-bundle power on its moduli space. G. Masbaum asks to study the specified action of the finite group \(E(\Sigma)\) on \(H_k\), determine its decomposition according to characters, and relate that decomposition to those of \(V(\Sigma)\) in the TQFT derived from \(U_q(\mathfrak{sl}_2)\) at a \((k+N)\)-th root of unity.

For a genuine linear complex action \(\rho_k\) on a finite-dimensional \(H_k\), character decomposition means the isotypic splitting
\[
 H_k\cong\bigoplus_{\chi\in\operatorname{Irr}(E(\Sigma))}
 W_\chi\otimes\C^{m_\chi(k)},\qquad
 m_\chi(k)=\frac1{|E(\Sigma)|}\sum_{g\in E(\Sigma)}
 \overline{\chi(g)}\operatorname{tr}(\rho_k(g)).
\]
The task is to identify the actual action, character multiplicities and comparison of summands; this general character identity alone is not their computation. No assumption that \(E(\Sigma)\) is abelian or has only one-dimensional characters is made.

The excerpt's ``this action'' lacks its defining antecedent. The group, moduli space, line bundle, action/lift, level range and parameter \(N\) are not defined. If the action is initially projective, the linear or central-extension convention must be specified before applying the displayed splitting. Nor is an identification between \(H_k\) and the TQFT vector space supplied. The exact \((k+N)\)-th-root condition is retained without substituting a guessed level shift or root order.
\subsection{Short English statement}
Compute the character decomposition of the specified finite-group action on the holomorphic section space, then compare it with decompositions in the stated \(\mathfrak{sl}_2\) TQFT. The action and geometric/level conventions are missing.
\subsection{Sources}
\begin{itemize}
\item Local UnsolvedMath registry, permanent ID 10400155, source code catalogue entry; statement attribution G. Masbaum; proposer field and proposal year absent; updated 2026-08-06.
\item Ohtsuki, \emph{Problems on invariants of knots and 3-manifolds} (2002), Problem 8.9, PDF page 125: \url{https://msp.org/gtm/2002/04/gtm-2002-04-024s.pdf}. Registry access date 2026-07-29; PDF font-structure extraction. Provenance distinguishes 231 numbered bold headings. Present primary-literature status and redistribution terms require review.
\item The supplied summary repeats the finite-group action question without precise additional definitions, references, or a verified resolution.
\end{itemize}
\end{document}
